\section{Verification purpose}

\textbf{Case ID:} \texttt{\CaseID}. This case verifies slope enhancement and aspect handling across two opposing valley sides.  Zero wind isolates the terrain contribution.  The verification question is whether the maximum regional spread rates on the 5-degree east side and 15-degree west side are each within 10 percent of their independent targets.

\section{Mathematical formulation and numerical implementation}

The surface spread calculation multiplies the no-wind/no-slope response by a slope-dependent factor and orients the heading direction upslope.  The slope raster supplies magnitude in degrees and the aspect raster supplies downslope azimuth.  Opposing aspects must therefore drive the two head fires away from the central valley ignition.

\section{Derivation of expected behavior}

For the configured Fuel Model 3 and moisture, the reference calculation gives \(1.92~\mathrm{m\,min^{-1}}\) on the 5-degree side and \(5.37~\mathrm{m\,min^{-1}}\) on the 15-degree side.  The steeper slope must have the larger head rate; this ordering is an additional qualitative invariant.

\section{Verification metrics and acceptance criteria}

For each half-domain, the metric is the maximum finite \texttt{vs} value. Regional maxima isolate the intended slope without averaging across the slope discontinuity or unburned cells. Each target uses 10 percent relative tolerance; both regional decisions are required for PASS.

\section{Simulation configuration and predicted outcome}

The 30 m grid has Fuel Model 3, zero wind, a 5-degree eastern slope with downslope aspect 270 degrees, a 15-degree western slope with downslope aspect 90 degrees, and central ignition. Predicted spread is outward and faster on the western, steeper side.

\section{Actual simulation results}

The required evidence is generated only from case-local ELMFIRE outputs. The current postprocessing status is \textbf{\VerificationStatus}; \PassingMetricCount{} of \MetricCount{} required metric rows pass across \VariantCount{} variant(s). A missing output is reported as NOT RUN or NOT EVALUABLE and is not interpreted as verification evidence.


\begin{figure}[H]
  \centering
  \IfFileExists{../figures/whole_domain_result.pdf}{%
    \includegraphics[width=0.82\textwidth,height=0.58\textheight,keepaspectratio]{whole_domain_result.pdf}%
  }{\fbox{\parbox{0.85\textwidth}{No case-local whole-domain ELMFIRE raster is available.}}}
  \caption{Representative whole-domain case-local ELMFIRE output. The generated figure records its source path, coordinate reference system, and variant.}
\end{figure}


\begin{figure}[H]
  \centering
  \IfFileExists{../figures/verification_summary.pdf}{%
    \includegraphics[width=\textwidth,height=0.58\textheight,keepaspectratio]{verification_summary.pdf}%
  }{\fbox{\parbox{0.85\textwidth}{No case-local summary figure is available.}}}
  \caption{Case-local expected and measured verification metrics. A missing measured value denotes unavailable or incomplete ELMFIRE output.}
\end{figure}

\section{Calculated metrics and verification decision}

\begin{center}
\begin{tabular}{@{}llll@{}}
\toprule
Variant / metric & Expected & Calculated & Status \\
\midrule
\MetricRows
\bottomrule
\end{tabular}
\end{center}

The authoritative overall decision is \textbf{\VerificationStatus} and is generated by \texttt{scripts/postprocess.py} from \texttt{outputs/metrics.json}; it is not recomputed in this report.

\section{Technical interpretation}

Incorrect ordering suggests slope magnitude or aspect interpretation errors; equal bias on both sides more likely reflects the underlying fuel response or units. No scientific conclusion should be drawn when the status is NOT RUN or NOT EVALUABLE.

\section{Scope and limitations}

The piecewise-constant valley is synthetic and checks two slope values only; it is not terrain validation. This case provides implementation-verification evidence only.
